\section{The one idea that explains everything}

Rusty Pomodoro is a Pomodoro timer. That part is easy. The interesting part is that it has
\emph{two entirely different user interfaces} --- a hand-built AppKit one and an egui
one --- and the timer logic is shared byte-for-byte between them.

That only works because of a strict layering rule, which is stated in the first line of
\file{src/mac\_shell.rs}:

\begin{lstlisting}[caption={The module's own doc comment states the contract.}]
//! Renderer-free macOS UI. Rust owns timer/storage; AppKit owns controls and drawing.
\end{lstlisting}

\subsection{The three layers}

\begin{figure}[H]
\centering
\begin{tikzpicture}[
  column/.style={rounded corners=3pt, line width=1pt, inner sep=5pt, align=center,
                 text width=33mm, minimum height=70mm, font=\scriptsize\sffamily},
]
  % ------------------------------------------------------------------ core
  \node[column, fill=corefill, draw=pastel!90!black] (core) at (0,0) {%
    \hd{\color{pastel!45!black}THE CORE}\\[2mm]
    \texttt{app.rs}\\[1.2mm] \texttt{timer.rs}\\[1.2mm]
    \texttt{config.rs}\\[1.2mm] \texttt{stats.rs}\\[3mm]
    Knows: what time it is, what activity comes\\next, which settings are valid,
    today's\\totals.\\[2.5mm]
    {\itshape No windows. No graphics.\\No AppKit. No egui.}};

  % ------------------------------------------------------------------ shell
  \node[column, fill=shellfill, draw=mint!80!black, right=27mm of core] (shell) {%
    \hd{\color{mint!60!black}THE SHELL}\\[2mm]
    \texttt{mac\_shell.rs}\\[1.2mm] \texttt{egui\_shell.rs}\\[1.2mm]
    \texttt{native.rs}\\[1.2mm] \texttt{hotkeys.rs}\\[1.2mm] \texttt{theme.rs}\\[3mm]
    Knows: window creation, event\\sources, repaint scheduling, rendering,
    menu\\bar, sound playback.\\[2.5mm]
    {\itshape Replaceable.}};

  % ------------------------------------------------------------------ os
  \node[column, fill=soft, draw=softink!55, dashed, right=27mm of shell] (os) {%
    \hd{\color{softink}THE PLATFORM}\\[2mm]
    \textbf{AppKit}\\[1.2mm]
    NSWindow, NSButton, NSSound,\\NSStatusBar, NSTimer\\[2.5mm]
    \hd{\color{softink}egui}\\[1.2mm]
    glow, winit, x11\\[2.5mm]
    {\itshape Differs per build.}};

  % ------------------------------------------------- core -> shell (intent)
  \draw[er, line width=1.5pt]
    (core.east) -- node[vlabel, midway, above=1.5mm, align=center, text=accent!60!black]
      {\textbf{intent flags}\\\texttt{hide\_requested}\\\texttt{notify\_requested}}
    ($(shell.west)+(0,9mm)$);

  % ------------------------------------------------- shell -> core (events)
  \draw[es, line width=1.5pt]
    ($(shell.west)+(0,-9mm)$) -- node[vlabel, midway, below=1.5mm, align=center, text=mint!60!black]
      {\textbf{callbacks}\\\texttt{tick()} \texttt{space}\\\texttt{willSleep}}
    (core.east);

  % ------------------------------------------------- shell -> platform
  \draw[e, line width=1.5pt]
    (shell.east) -- node[vlabel, midway, above=1.5mm, align=center]
      {\textbf{effect calls}\\\texttt{NSWindow}, \texttt{NSSound},\\\texttt{orderOut}, \texttt{save}}
    ($(os.west)+(0,9mm)$);
  \draw[e, line width=1.5pt]
    ($(os.west)+(0,-9mm)$) -- node[vlabel, midway, below=1.5mm, align=center]
      {\textbf{effect calls}\\\texttt{egui::Context},\\\texttt{ViewportCommand}, \texttt{save}}
    (shell.east);

  % ------------------------------------------------- group brackets
  \draw[decorate, decoration={brace, amplitude=4pt, mirror}, thick, pastel!45!black]
    ($(core.north west)+(-3mm,2.5mm)$) -- ($(core.north east)+(3mm,2.5mm)$)
    node[midway, above=2mm, font=\tiny\sffamily\bfseries, text=pastel!45!black]
      {shared by every build, unit-tested};

  \draw[decorate, decoration={brace, amplitude=4pt, mirror}, thick, mint!80!black]
    ($(shell.north west)+(-3mm,2.5mm)$) -- ($(shell.north east)+(3mm,2.5mm)$)
    node[midway, above=2mm, font=\tiny\sffamily\bfseries, text=mint!60!black]
      {swapped per target OS};

  \node[note, below=6mm of core.south, text width=33mm, font=\tiny\sffamily]
    {The core could be compiled with zero\\platform dependencies; \texttt{cfg}
     gates keep that true.};
\end{tikzpicture}
\caption{The layering. The core never calls a platform API; it only raises flags. The
shell reads them, performs the effect, and clears them with \texttt{std::mem::take}.}
\label{fig:layers}
\end{figure}

\subsection{The mechanism: five boolean fields}

The whole boundary is five public booleans on \file{App}, from \file{src/app.rs} lines 22--29:

\begin{lstlisting}[caption={Every side effect the core wants is a public boolean field.}]
pub hide_requested:   bool,   // shell should order the window out
pub show_requested:   bool,   // shell should bring the window forward
pub notify_requested: bool,   // shell should play the completion sound
pub tick_requested:   bool,   // shell should play the per-second tick sound
pub exit_requested:   bool,   // shell should terminate the process
\end{lstlisting}

The core never calls \texttt{NSBeep()} or \texttt{request\_repaint}. It only ever writes
\texttt{true}. The shell then calls \texttt{std::mem::take(\&mut app.notify\_requested)},
which reads the flag \emph{and} resets it to \texttt{false} in one step:

\begin{lstlisting}[caption={Both shells drain the flags identically, at the end of a frame.}]
// src/mac_shell.rs, refresh()
if std::mem::take(&mut app.hide_requested)   { views.window.orderOut(None); }
let show = std::mem::take(&mut app.show_requested);
if std::mem::take(&mut app.notify_requested) { play(app.settings.notification_sound); }
if std::mem::take(&mut app.tick_requested)   { play(Sound::Tink); }
\end{lstlisting}

Because the flag is cleared by the same read that tests it, an intent can never fire
twice. And because the core only writes booleans, it never needs to know whether a
shell exists at all --- the same \file{App} compiles inside both shells unmodified.

\begin{figure}[H]
\centering
\begin{tikzpicture}[
  fw/.style={rounded corners=2pt,draw=ink!50,fill=white,align=center,
             inner sep=5pt,font=\scriptsize\sffamily,minimum height=10mm},
  key/.style={rounded corners=2pt,draw=accent,fill=accent!14,align=center,
              inner sep=5pt,line width=1.2pt,font=\scriptsize\sffamily\bfseries,
              text=accent!60!black,minimum height=10mm},
]
  \node[fw] (n1) {\textbf{core}\\ decides something\\ should happen};
  \node[key, right=6mm of n1] (n2) {sets a flag\\\texttt{= true}};
  \node[fw, right=6mm of n2] (n3) {\textbf{shell}\\ reaches end of\\its frame callback};
  \node[key, right=6mm of n3] (n4) {\texttt{mem::take}\\ returns true\\and clears};
  \node[fw, right=6mm of n4] (n5) {\textbf{shell}\\ performs the\\real effect};
  \foreach \a/\b in {n1/n2,n2/n3,n3/n4,n4/n5} \draw[e] (\a) -- (\b);

  \draw[ethin, dashed, accent!70] (n2.south) ++(0,-1mm) -- ++(0,-6mm)
        -| node[vlabel, below, pos=0.78, text=accent!60!black]
        {\emph{nothing else: no callback, no closure, no channel}}
        ($(n5.south)+(0,-6mm)$);

  \node[note, below=15mm of n3.5, text width=145mm] {%
    The core can be compiled, unit-tested, and reasoned about with no UI at all.%
    It also means a second UI does not add a single line of new logic --- it
    only adds another consumer for the same five flags.};
\end{tikzpicture}
\caption{The intent-flag protocol, as a five-step cycle.}
\label{fig:flags}
\end{figure}

\subsection{What one second of running looks like}

\begin{figure}[H]
\centering
\begin{tikzpicture}[
  fw/.style={rounded corners=2pt,draw=ink!50,fill=white,align=center,
             inner sep=4pt,font=\scriptsize\sffamily,minimum height=11mm,text width=15mm},
  key/.style={fw, draw=accent, fill=accent!14, line width=1.2pt,
              text=accent!60!black, font=\scriptsize\sffamily\bfseries},
]
  \node[fw] (n1) {event fires\\\texttt{NSTimer}\\\scriptsize or \texttt{logic()}};
  \node[key, right=4mm of n1] (n2) {\texttt{app.tick()}\\\scriptsize reads the clock\\decides if the second\\changed};
  \node[fw, right=4mm of n2] (n3) {update labels\\\scriptsize \texttt{refresh()}\\\scriptsize or \texttt{effects()}};
  \node[fw, right=4mm of n3] (n4) {consume flags\\\scriptsize \texttt{mem::take}};
  \node[fw, right=4mm of n4] (n5) {save settings\\\scriptsize if dirty};
  \node[fw, right=4mm of n5] (n6) {schedule the\\\scriptsize next wakeup};
  \foreach \a/\b in {n1/n2,n2/n3,n3/n4,n4/n5,n5/n6} \draw[e] (\a) -- (\b);

  \draw[ethin, dashed, accent!70] (n2.north) ++(0,1mm) -- ++(0,7mm)
        -| node[vlabel, above, pos=0.80, text=accent!60!black]
        {\emph{only when the displayed second actually changed}}
        ($(n6.north)+(0,7mm)$);

  \node[note, below=8mm of n3.5, text width=150mm] {%
    No thread, no async runtime, no polling loop. When the timer is \emph{Idle},
    step six schedules nothing at all, so the process does literally zero work per second.};
\end{tikzpicture}
\caption{The per-second cycle in the AppKit shell. \texttt{refresh()} ends by arming
exactly one one-shot \texttt{NSTimer} for the instant the next displayed digit changes.}
\label{fig:second}
\end{figure}